\documentclass{article}
\usepackage{geometry}
\geometry{a4paper,scale=0.8}
\usepackage[utf8]{inputenc}

\usepackage{amsmath, amssymb}
\usepackage{amsthm}
\usepackage{enumitem}
\usepackage{tikz-cd}
\usepackage{mathrsfs}

\newtheoremstyle{breakstyle}
    {5pt}
    {10pt}
    % {\normalfont}
    {\itshape}
    {0pt}
    {\bfseries}
    {\newline\indent}
    {0pt}
    {\thmname{#1}\thmnumber{ #2}\thmnote{ \textbf{[#3]}}}
\theoremstyle{breakstyle}
\newtheorem{exercise}{Exercise}[section]
\renewcommand{\theexercise}{\thesection.\Alph{exercise}}
\newtheorem{theorem}{Theorem}[section]
\newtheorem{proposition}{Proposition}[section]
\newtheorem{corollary}{Corollary}[section]
\newtheorem{remark}{Remark}[section]
\newtheorem{definition}{Definition}[section]
\newtheorem{lemma}{Lemma}[section]

\DeclareMathOperator{\id}{Id}
\DeclareMathOperator{\mor}{Mor}
\DeclareMathOperator{\coker}{Coker}
\let\ker\relax
\DeclareMathOperator{\ker}{Ker}
\DeclareMathOperator{\im}{Im}
\let\lim\relax
\DeclareMathOperator{\lim}{Lim}
\DeclareMathOperator{\colim}{Colim}
\let\hom\relax
\DeclareMathOperator{\hom}{Hom}
\let\mod\relax
\DeclareMathOperator{\mod}{Mod}
\DeclareMathOperator{\pr}{pr}
\DeclareMathOperator{\bil}{Bil}
\let\dim\relax
\DeclareMathOperator{\dim}{Dim}
\let\gcd\relax
\DeclareMathOperator{\gcd}{Gcd}

\title{Chapter 2 - Sheaves}
\author{Flandre}
\date{\today}

\begin{document}
\maketitle

\section{Motivating example: The sheaf of smooth functions}

\begin{exercise}
Show that this is the only maximal ideal of $O_{p}$ . (Hint: show that every
element of $O_{p} \setminus m_{p}$ is invertible.)

\begin{proof}
    For $\forall \left( f,U \right) \in \mathscr{O}_{p}\backslash
    \mathfrak{m}_{p}$, we have
    $$
    \left( \frac{1}{f}, U'\right)\in \mathscr{O}_{p}\backslash \mathfrak{m}_{p}
    $$
    in a neighbourhood $U'$ of $p$, hence invertible.
    \qed
\end{proof}

\end{exercise}

\begin{exercise}[Skipped. I haven't learnt Differential Geometry yet]
EXERCISE FOR THOSE WITH DIFFERENTIAL GEOMETRIC BACKGROUND.

Prove this. (Rhetorical question for experts: what goes wrong if the sheaf of
continuous functions is substituted for the sheaf of smooth functions? What goes
wrong if you use the sheaf of $C^{1}$ functions?)
\end{exercise}

\section{Definition of sheaf and presheaf}

\begin{exercise}[Trivial]
EXERCISE FOR CATEGORY-LOVERS: “A PRESHEAF IS THE SAME AS A CONTRAVARIANT
FUNCTOR”.

Given any topological space X, we have a “category of open sets” (Example
1.1.9), where the objects are the open sets and the morphisms are inclusions.
Verify that the data of a presheaf is precisely the data of a contravariant
functor from the category of open sets of X to the category of sets. (This
interpretation is surprisingly useful.)  
\end{exercise}

\end{document}
